% ═══════════════════════════════════════════════════════════════════════════
% Chapter 13: Data Collection Methodology
% ═══════════════════════════════════════════════════════════════════════════

\chapter{Data Collection Methodology} \label{chap:methodology}

\section{Overview}

This chapter details the experimental protocols, subject recruitment criteria, session procedures, and data management practices employed in collecting the validation dataset comprising 55 sessions and 622 annotated repetitions.

\section{Subject Recruitment}

\subsection{Inclusion Criteria}

Participants were required to meet the following criteria:
\begin{itemize}
    \item Age 18-45 years
    \item Minimum 2 years of resistance training experience
    \item Ability to perform all six compound exercises with proper technique
    \item No current musculoskeletal injuries limiting barbell exercises
    \item Signed informed consent form
\end{itemize}

\subsection{Demographics}

\begin{table}[h]
    \centering
    \caption{Subject demographics (n=12)}
    \label{tab:demographics}
    \begin{tabular}{lc}
        \toprule
        \textbf{Characteristic} & \textbf{Mean ± SD} \\
        \midrule
        Age (years) & 28.4 ± 5.2 \\
        Height (cm) & 176.3 ± 8.1 \\
        Body mass (kg) & 78.5 ± 12.3 \\
        Training experience (years) & 4.8 ± 2.1 \\
        \midrule
        \textbf{Sex Distribution} & \\
        Male & 8 (67\%) \\
        Female & 4 (33\%) \\
        \bottomrule
    \end{tabular}
\end{table}

\section{Exercise Selection}

Six compound barbell exercises were selected based on:
\begin{enumerate}
    \item Prevalence in strength training programs
    \item Varied movement patterns (hinge, squat, push, pull, Olympic lift)
    \item Different velocity profiles and ranges of motion
\end{enumerate}

\subsection{Exercise Descriptions}

\begin{description}
    \item[Back Squat:] Barbell positioned on upper back, squat to parallel or below, drive through heels to lockout.
    \item[Bench Press:] Supine on bench, barbell lowered to mid-chest, press to full elbow extension.
    \item[Deadlift:] Barbell on floor, hip hinge to grip, extend hips and knees to lockout.
    \item[Snatch:] Barbell from floor to overhead in one motion, rapid pull and turnover.
    \item[Clean:] Barbell from floor to shoulder position, rapid pull and front rack reception.
    \item[Barbell Row:] Bent-over position, pull barbell to lower abdomen, controlled descent.
\end{description}

\section{Session Protocol}

\subsection{Pre-Session Procedures}

\begin{enumerate}
    \item \textbf{Consent verification:} Confirm informed consent on file.
    \item \textbf{Anthropometric measurement:} Height, body mass recorded.
    \item \textbf{Warm-up:} 5-10 minutes light cardio, dynamic stretching.
    \item \textbf{Equipment setup:} Attach IMU node to barbell collar, calibrate.
    \item \textbf{Camera calibration:} Position RealSense D455, verify FRAME\_SYNC.
\end{enumerate}

\subsection{Loading Scheme}

Approximate 1RM estimated from submaximal lifts (~70\% load for 8-10 reps). Working sets performed at 65-85\% estimated 1RM:

\begin{table}[h]
    \centering
    \caption{Loading scheme by exercise}
    \label{tab:loading}
    \begin{tabular}{lccc}
        \toprule
        \textbf{Exercise} & \textbf{First Set (kg)} & \textbf{Working \%1RM} & \textbf{Target Reps} \\
        \midrule
        Back Squat & 40-60 & 70-80\% & 5-8 \\
        Bench Press & 20-40 & 70-80\% & 5-8 \\
        Deadlift & 60-80 & 75-85\% & 3-5 \\
        Snatch & 20-30 & 65-75\% & 3-5 \\
        Clean & 30-40 & 70-80\% & 3-5 \\
        Barbell Row & 30-40 & 70-80\% & 6-10 \\
        \bottomrule
    \end{tabular}
\end{table}

\subsection{Rest Intervals}

\begin{itemize}
    \item \textbf{Warm-up sets:} 60-90 seconds rest
    \item \textbf{Working sets:} 180-240 seconds rest
    \item \textbf{Between exercises:} 300+ seconds rest
\end{itemize}

\section{IMU Calibration Procedure}

\subsection{Static Calibration}

Prior to each session, the barbell is placed on the floor and remains motionless for 5 seconds. The IMU accumulates samples to compute:

\begin{align}
    \vec{b}_g &= \frac{1}{N}\sum_{i=1}^{N}\vec{\omega}_i \quad \text{(gyro bias)} \\
    s_a &= \frac{1g}{\frac{1}{N}\sum_{i=1}^{N}|\vec{a}_i|} \quad \text{(accel scale)}
\end{align}

Calibration validation requires $\sigma_a < 0.02g$ during the static period.

\subsection{Dynamic Validation}

After calibration, the operator performs 2-3 submaximal reps to verify:
\begin{itemize}
    \item Data transmission intact (ESP-NOW ACK >99\%)
    \item Velocity traces plausible (0-2 m/s range for typical lifts)
    \item No dropped samples or timestamp anomalies
\end{itemize}

\section{Camera Ground Truth Collection}

\subsection{Mounting Position}

The RealSense D455 is mounted on a tripod, positioned 3-4 meters from the barbell path, at approximately hip height. The camera is perpendicular to the plane of motion to minimize perspective distortion.

\subsection{Synchronization Verification}

At session start, a simultaneous clap or sharp acceleration is performed. The IMU accelerometer spike and camera frame sync signal are cross-correlated to determine the temporal offset, which is applied to align all subsequent samples.

\section{Annotation Procedure}

\subsection{Manual Labeling}

After each session, video and IMU traces are reviewed frame-by-frame to annotate rep boundaries. The annotation interface displays:
\begin{itemize}
    \item Synchronized video playback
    \item IMU acceleration/gyro traces
    \item Computed velocity and position
    \item Drag handles for boundary adjustment
\end{itemize}

\subsection{Quality Control}

\begin{enumerate}
    \item \textbf{First pass:} Auto-segmentation provides candidate boundaries.
    \item \textbf{Second pass:} Human annotator adjusts boundaries to match video evidence.
    \item \textbf{Third pass:} Senior annotator reviews all annotations with F1 <0.95 vs. auto-segmentation.
\end{enumerate}

\section{Data Management}

\subsection{File Organization}

Each session is stored in a schema-versioned directory:

\begin{verbatim}
sessions/session_YYYYMMDD_HHMMSS/
├── metadata.json
├── manifest.json
├── imu/
│   └── raw_imu.csv
├── video/
│   ├── color.mp4
│   └── depth.npy
└── annotations/
    └── rep_segments.json
\end{verbatim}

\subsection{Metadata Fields}

Over 120 metadata fields captured including:
\begin{itemize}
    \item Subject demographics
    \item Training context (phase, set number, load, RPE)
    \item Equipment specifications
    \item Environmental conditions (gym location, temperature)
    \item Session quality flags
\end{itemize}

\section{Ethical Considerations}

\subsection{Informed Consent}

All participants signed informed consent forms detailing:
\begin{itemize}
    \item Purpose of data collection
    \item Risks (minimal: barbell exercise)
    \item Data usage rights
    \item Right to withdraw at any time
\end{itemize}

\subsection{Data Anonymization}

Subject identifiers are replaced with coded IDs. Original consent forms linking codes to names are stored separately, accessible only to principal investigators.

\section{Data Quality Assurance}

\subsection{Automated Checks}

Upon session completion, automated validation scripts verify:
\begin{itemize}
    \item No gaps >50ms in IMU stream
    \item Acceleration magnitudes within $\pm$16 g
    \item Gyro rates within $\pm$2000 dps
    \item ESP-NOW sequence numbers contiguous
\end{itemize}

\subsection{Manual Review}

Sessions with anomalies flagged for expert review before inclusion in the research dataset.

\section{Limitations}

\begin{itemize}
    \item \textbf{Sample size:} 12 subjects limits generalizability to diverse populations.
    \item \textbf{Exercise selection:} Only six exercises studied; unilateral movements not evaluated.
    \item \textbf{Loading range:} 65-85\% 1RM; heavier (>90\%) and lighter (<60\%) loads underrepresented.
    \item \textbf{Fatigue state:} Sessions conducted fresh; no data on velocity under cumulative fatigue.
\end{itemize}

\section{Summary}

This chapter detailed the systematic methodology for collecting the validation dataset. Rigid protocols for subject recruitment, calibration, synchronization, and annotation ensure data quality suitable for research-grade algorithm development and validation.

\newpage